\section{Cancer as Impedance Decoupling}
\label{sec:cancer_impedance}

\subsection{Hypothesis}

Normal tissue maintains cooperative standing-wave patterns through
impedance-matched cell--cell coupling via gap junctions.  Each cell
presents a characteristic impedance $Z_{\mathrm{cell}}$ determined
by its membrane composition, ion-channel distribution, and
cytoskeletal stiffness.  When these impedances are matched, the
morphogenetic standing wave---the spatially-periodic field that
encodes tissue identity and growth regulation---propagates freely
through the multicellular network with minimal reflection
($\Gamma \approx 0$).

Cancer, in this framework, is \emph{impedance decoupling}.  A
mutation that alters membrane lipid composition, channel density,
or cytoskeletal tension shifts $Z_{\mathrm{cell}}$ away from the
tissue-matched value.  The reflection coefficient at the
mutant/healthy boundary rises:
\begin{equation}
  \Gamma_{\mathrm{tumour}} 
  = \frac{Z_{\mathrm{mutant}} - Z_{\mathrm{healthy}}}
         {Z_{\mathrm{mutant}} + Z_{\mathrm{healthy}}}.
\end{equation}
When $|\Gamma|$ exceeds a critical threshold, the cell is
effectively isolated from the tissue standing wave.  The ``stop
dividing'' signal---a morphogenetic cavity mode---can no longer
reach it.  Unregulated proliferation follows.

\subsection{Metastasis}

Metastasis acquires a natural interpretation: a cell whose
$Z_{\mathrm{cell}}$ has drifted far enough from its tissue of
origin may, by chance, impedance-match a \emph{different} tissue
type.  Migration to that tissue restores $\Gamma \approx 0$ for
the new environment, enabling colonisation.

\subsection{Tumour Suppression}

The p53 tumour suppressor protein functions, in this model, as an
impedance calibration checkpoint.  When a cell's impedance drifts
beyond the tissue-matched tolerance band, p53 triggers apoptosis---a
controlled network reset that removes the mismatched element before
the standing wave can collapse.

\subsection{Testable Predictions}

\begin{enumerate}
  \item Tumour impedance spectra, measured via Electrical Impedance
        Spectroscopy (EIS), should show systematic loss of the
        standing-wave resonance peaks present in healthy tissue.
  \item The tumour/healthy boundary should exhibit elevated
        $|\Gamma|$ at the morphogenetic standing-wave frequency.
  \item Cell lines with p53 knockout should show larger impedance
        drift variance than wild-type, measurable via broadband EIS.
\end{enumerate}

\subsection{Research Avenues for Treatment}

Four treatment strategies follow directly from the impedance model:

\begin{description}
  \item[Impedance-targeted RF ablation.]
    Design radiofrequency fields whose frequency and polarisation
    are impedance-matched to the tumour ($\Gamma_{\mathrm{tumour}}
    \to 0$, maximum energy absorption) while reflecting off healthy
    tissue ($\Gamma_{\mathrm{healthy}}$ large).  This is
    frequency-selective ablation guided by pre-operative EIS mapping.

  \item[Morphogenetic recoupling.]
    Apply external fields at the healthy tissue's standing-wave
    frequency to re-establish the cooperative mode.  Cells that can
    recouple differentiate; cells too far drifted to recouple are
    driven to apoptosis by the impedance mismatch stress.

  \item[EIS-guided diagnostics.]
    Tumour boundaries are impedance discontinuities.  The reflection
    coefficient $\Gamma$ at the tumour/healthy interface is large
    and spatially mappable via multi-electrode EIS arrays, providing
    sub-millimetre margin delineation without ionising radiation.

  \item[Impedance-matched drug targeting.]
    Design drug molecules whose topological impedance
    $Z_{\mathrm{topo}}$ matches the tumour cell membrane, giving
    $\Gamma \approx 0$ for selective uptake.  For healthy cells with
    a different $Z$, $\Gamma$ is large and the drug is reflected.
\end{description}
